% ---------------------------------------------------------
% TABELA EMENTA 25: Geografia — Ano 1
% ---------------------------------------------------------
\begin{xltabular}{\linewidth}{|X|p{2.5cm}|p{2.5cm}|}
\hline
\multirow{3}{=}{\textcolor{ifscgreen}{\textbf{Unidade Curricular:}}\\[3pt]\textbf{\large Geografia — Ano 1}} & \multicolumn{2}{p{\dimexpr5cm+2\tabcolsep+\arrayrulewidth\relax}|}{\textcolor{ifscgreen}{\textbf{Semestre:}} \textbf{1º e 2º}} \\ \cline{2-3}
 & \textcolor{ifscgreen}{\textbf{CH EaD*:}} & \textcolor{ifscgreen}{\textbf{CH Total*:}} \\ \cline{2-3}
 & \textbf{00 h} & \textbf{80 h} \\ \hline
\endhead
\multicolumn{3}{|p{\dimexpr\linewidth-2\tabcolsep-2\arrayrulewidth\relax}|}{\textcolor{ifscgreen}{\textbf{Objetivos:}}} \\ \hline
\multicolumn{3}{|p{\dimexpr\linewidth-2\tabcolsep-2\arrayrulewidth\relax}|}{
\begin{itemize}[leftmargin=*,noitemsep,topsep=2pt]
 \item \textit{[Objetivos de aprendizagem pendentes de preenchimento pelo docente responsável]}
\end{itemize}
} \\ \hline
\multicolumn{3}{|p{\dimexpr\linewidth-2\tabcolsep-2\arrayrulewidth\relax}|}{\textcolor{ifscgreen}{\textbf{Conteúdos:}}} \\ \hline
\multicolumn{3}{|p{\dimexpr\linewidth-2\tabcolsep-2\arrayrulewidth\relax}|}{
\begin{itemize}[leftmargin=*,noitemsep,topsep=2pt]
 \item \textit{[Conteúdo programático pendente de preenchimento pelo docente responsável]}
\end{itemize}
} \\ \hline
\multicolumn{3}{|p{\dimexpr\linewidth-2\tabcolsep-2\arrayrulewidth\relax}|}{\textcolor{ifscgreen}{\textbf{Estratégias de Ensino e Aprendizagem:}}} \\ \hline
\multicolumn{3}{|p{\dimexpr\linewidth-2\tabcolsep-2\arrayrulewidth\relax}|}{\textit{[Metodologia de ensino e avaliação pendente de preenchimento pelo docente responsável]}} \\ \hline
\multicolumn{3}{|p{\dimexpr\linewidth-2\tabcolsep-2\arrayrulewidth\relax}|}{\textcolor{ifscgreen}{\textbf{Bibliografia Básica:}}} \\ \hline
\multicolumn{3}{|p{\dimexpr\linewidth-2\tabcolsep-2\arrayrulewidth\relax}|}{1. \textit{[1º título da Bibliografia Básica pendente de indicação docente e validação no Parecer da Biblioteca do Câmpus]}
2. \textit{[2º título da Bibliografia Básica pendente de indicação docente e validação no Parecer da Biblioteca do Câmpus]}} \\ \hline
\multicolumn{3}{|p{\dimexpr\linewidth-2\tabcolsep-2\arrayrulewidth\relax}|}{\textcolor{ifscgreen}{\textbf{Bibliografia Complementar:}}} \\ \hline
\multicolumn{3}{|p{\dimexpr\linewidth-2\tabcolsep-2\arrayrulewidth\relax}|}{1. \textit{[1º título da Bibliografia Complementar pendente de indicação docente e validação no Parecer da Biblioteca do Câmpus]}
2. \textit{[2º título da Bibliografia Complementar pendente de indicação docente e validação no Parecer da Biblioteca do Câmpus]}
3. \textit{[3º título da Bibliografia Complementar pendente de indicação docente e validação no Parecer da Biblioteca do Câmpus]}} \\ \hline
\end{xltabular}

\vspace{0.4cm}


% ---------------------------------------------------------
% TABELA EMENTA 26: Geografia — Ano 3
% ---------------------------------------------------------
\begin{xltabular}{\linewidth}{|X|p{2.5cm}|p{2.5cm}|}
\hline
\multirow{3}{=}{\textcolor{ifscgreen}{\textbf{Unidade Curricular:}}\\[3pt]\textbf{\large Geografia — Ano 3}} & \multicolumn{2}{p{\dimexpr5cm+2\tabcolsep+\arrayrulewidth\relax}|}{\textcolor{ifscgreen}{\textbf{Semestre:}} \textbf{5º e 6º}} \\ \cline{2-3}
 & \textcolor{ifscgreen}{\textbf{CH EaD*:}} & \textcolor{ifscgreen}{\textbf{CH Total*:}} \\ \cline{2-3}
 & \textbf{00 h} & \textbf{80 h} \\ \hline
\endhead
\multicolumn{3}{|p{\dimexpr\linewidth-2\tabcolsep-2\arrayrulewidth\relax}|}{\textcolor{ifscgreen}{\textbf{Objetivos:}}} \\ \hline
\multicolumn{3}{|p{\dimexpr\linewidth-2\tabcolsep-2\arrayrulewidth\relax}|}{
\begin{itemize}[leftmargin=*,noitemsep,topsep=2pt]
 \item \textit{[Objetivos de aprendizagem pendentes de preenchimento pelo docente responsável]}
\end{itemize}
} \\ \hline
\multicolumn{3}{|p{\dimexpr\linewidth-2\tabcolsep-2\arrayrulewidth\relax}|}{\textcolor{ifscgreen}{\textbf{Conteúdos:}}} \\ \hline
\multicolumn{3}{|p{\dimexpr\linewidth-2\tabcolsep-2\arrayrulewidth\relax}|}{
\begin{itemize}[leftmargin=*,noitemsep,topsep=2pt]
 \item \textit{[Conteúdo programático pendente de preenchimento pelo docente responsável]}
\end{itemize}
} \\ \hline
\multicolumn{3}{|p{\dimexpr\linewidth-2\tabcolsep-2\arrayrulewidth\relax}|}{\textcolor{ifscgreen}{\textbf{Estratégias de Ensino e Aprendizagem:}}} \\ \hline
\multicolumn{3}{|p{\dimexpr\linewidth-2\tabcolsep-2\arrayrulewidth\relax}|}{\textit{[Metodologia de ensino e avaliação pendente de preenchimento pelo docente responsável]}} \\ \hline
\multicolumn{3}{|p{\dimexpr\linewidth-2\tabcolsep-2\arrayrulewidth\relax}|}{\textcolor{ifscgreen}{\textbf{Bibliografia Básica:}}} \\ \hline
\multicolumn{3}{|p{\dimexpr\linewidth-2\tabcolsep-2\arrayrulewidth\relax}|}{1. \textit{[1º título da Bibliografia Básica pendente de indicação docente e validação no Parecer da Biblioteca do Câmpus]}
2. \textit{[2º título da Bibliografia Básica pendente de indicação docente e validação no Parecer da Biblioteca do Câmpus]}} \\ \hline
\multicolumn{3}{|p{\dimexpr\linewidth-2\tabcolsep-2\arrayrulewidth\relax}|}{\textcolor{ifscgreen}{\textbf{Bibliografia Complementar:}}} \\ \hline
\multicolumn{3}{|p{\dimexpr\linewidth-2\tabcolsep-2\arrayrulewidth\relax}|}{1. \textit{[1º título da Bibliografia Complementar pendente de indicação docente e validação no Parecer da Biblioteca do Câmpus]}
2. \textit{[2º título da Bibliografia Complementar pendente de indicação docente e validação no Parecer da Biblioteca do Câmpus]}
3. \textit{[3º título da Bibliografia Complementar pendente de indicação docente e validação no Parecer da Biblioteca do Câmpus]}} \\ \hline
\end{xltabular}

\vspace{0.4cm}


% ---------------------------------------------------------
% TABELA EMENTA 27: História — Ano 2
% ---------------------------------------------------------
\begin{xltabular}{\linewidth}{|X|p{2.5cm}|p{2.5cm}|}
\hline
\multirow{3}{=}{\textcolor{ifscgreen}{\textbf{Unidade Curricular:}}\\[3pt]\textbf{\large História — Ano 2}} & \multicolumn{2}{p{\dimexpr5cm+2\tabcolsep+\arrayrulewidth\relax}|}{\textcolor{ifscgreen}{\textbf{Semestre:}} \textbf{3º e 4º (História I)}} \\ \cline{2-3}
 & \textcolor{ifscgreen}{\textbf{CH EaD*:}} & \textcolor{ifscgreen}{\textbf{CH Total*:}} \\ \cline{2-3}
 & \textbf{00 h} & \textbf{80 h} \\ \hline
\endhead
\multicolumn{3}{|p{\dimexpr\linewidth-2\tabcolsep-2\arrayrulewidth\relax}|}{\textcolor{ifscgreen}{\textbf{Objetivos:}}} \\ \hline
\multicolumn{3}{|p{\dimexpr\linewidth-2\tabcolsep-2\arrayrulewidth\relax}|}{
\begin{itemize}[leftmargin=*,noitemsep,topsep=2pt]
 \item Compreender a origem da humanidade e as primeiras levas migratórias e diaspóricas.
 \item Investigar as rotas migratórias dos primeiros grupos humanos para fora da África e os impactos dessas migrações nas diferentes regiões do mundo.
 \item Reconhecer a formação de comunidades ao redor do mundo e sua contribuição para a diversidade cultural global.
 \item Reconhecer os conhecimentos e tecnologias desenvolvidos por povos africanos e indígenas a partir das interações com o meio ambiente.
 \item Compreender a história e cultura dos povos indígenas da América.
 \item Analisar os impactos do contato europeu na vida e nas tradições dos povos indígenas, incluindo os processos de aculturação e resistência.
 \item Refletir sobre a resistência e a luta dos povos indígenas pela preservação de suas terras, culturas e direitos.
 \item Valorizar o intercâmbio cultural e a formação de comunidades diaspóricas.
 \item Analisar os processos de intercâmbio cultural entre diferentes comunidades diaspóricas, identificando as contribuições culturais, linguísticas, religiosas e artísticas dessas comunidades.
\end{itemize}
} \\ \hline
\multicolumn{3}{|p{\dimexpr\linewidth-2\tabcolsep-2\arrayrulewidth\relax}|}{\textcolor{ifscgreen}{\textbf{Conteúdos:}}} \\ \hline
\multicolumn{3}{|p{\dimexpr\linewidth-2\tabcolsep-2\arrayrulewidth\relax}|}{
\begin{itemize}[leftmargin=*,noitemsep,topsep=2pt]
 \item Origem da humanidade em África e as primeiras levas migratórias e diaspóricas: teorias e evidências científicas sobre a origem do Homo sapiens na África; rotas migratórias e suas consequências; formação de comunidades diaspóricas ao redor do mundo.
 \item Contribuições africanas para o desenvolvimento de conhecimentos e tecnologias: agricultura, metalurgia, medicina, matemática e astronomia.
 \item História e cultura dos povos indígenas da América: culturas, línguas, religiões e formas de organização social; impactos do contato europeu; resistência e luta pela preservação de terras, culturas e direitos.
 \item Intercâmbio cultural e formação de comunidades diaspóricas: contribuições culturais, linguísticas, religiosas e artísticas para a formação de novas identidades culturais.
\end{itemize}
} \\ \hline
\multicolumn{3}{|p{\dimexpr\linewidth-2\tabcolsep-2\arrayrulewidth\relax}|}{\textcolor{ifscgreen}{\textbf{Estratégias de Ensino e Aprendizagem:}}} \\ \hline
\multicolumn{3}{|p{\dimexpr\linewidth-2\tabcolsep-2\arrayrulewidth\relax}|}{- **Metodologia:** Abordagem participativa, crítica e contextualizada, amparada na pedagogia histórico-crítica, priorizando a relação horizontal entre professor e estudantes. Debates, análises de fontes históricas, atividades práticas e trabalhos em grupo, relacionando os conteúdos às experiências dos estudantes e à realidade social. Contextualização dos conteúdos históricos com acontecimentos e processos sociais atuais. Problematização das relações de poder e das desigualdades sociais na construção e interpretação da história.
- **Avaliação:** Não especificado no documento fonte.} \\ \hline
\multicolumn{3}{|p{\dimexpr\linewidth-2\tabcolsep-2\arrayrulewidth\relax}|}{\textcolor{ifscgreen}{\textbf{Bibliografia Básica:}}} \\ \hline
\multicolumn{3}{|p{\dimexpr\linewidth-2\tabcolsep-2\arrayrulewidth\relax}|}{1. BOULOS JÚNIOR, Alfredo. \textbf{360º: história sociedade \& cidadania}. São Paulo: FTD, 2015.
2. VICENTINO, Cláudio; DORIGO, Gianpaolo. \textbf{História geral e do Brasil}. São Paulo: Scipione, 2010.} \\ \hline
\multicolumn{3}{|p{\dimexpr\linewidth-2\tabcolsep-2\arrayrulewidth\relax}|}{\textcolor{ifscgreen}{\textbf{Bibliografia Complementar:}}} \\ \hline
\multicolumn{3}{|p{\dimexpr\linewidth-2\tabcolsep-2\arrayrulewidth\relax}|}{1. CERTEAU, Michel de. \textbf{A invenção do cotidiano: artes de fazer}. Petrópolis, RJ: Vozes, 2009.
2. RICKLEFS, Robert E. \textbf{A economia da natureza}. Rio de Janeiro: Guanabara Koogan, 2010.
3. CAREGNATO, Célia Elizabete; BOMBASSARO, Luiz Carlos (org.). \textbf{Diversidade cultural: viver diferenças e enfrentar desigualdades na educação}. Erechim, RS: Novello \& Carbonelli, 2013.
4. MIGNOLO, Walter. \textbf{Histórias Globais/projetos Locais}. Belo Horizonte: Editora UFMG, 2003.
5. GALEANO, Eduardo. \textbf{As caras e as máscaras}. Porto Alegre: L\&PM Editores, 2004.
6. HOLLANDA, Heloísa Buarque de (org.). \textbf{Pensamento Feminista Hoje: perspectivas decoloniais}. Rio de Janeiro: Bazar do Tempo, 2020.
7. SANTOS, Boaventura de Sousa; MARTINS, Bruno Sena (org.). \textbf{O pluriverso dos direitos humanos: a diversidade das lutas pela dignidade}. Belo Horizonte: Autêntica, 2019.} \\ \hline
\end{xltabular}

\vspace{0.4cm}


% ---------------------------------------------------------
% TABELA EMENTA 28: História — Ano 3
% ---------------------------------------------------------
\begin{xltabular}{\linewidth}{|X|p{2.5cm}|p{2.5cm}|}
\hline
\multirow{3}{=}{\textcolor{ifscgreen}{\textbf{Unidade Curricular:}}\\[3pt]\textbf{\large História — Ano 3}} & \multicolumn{2}{p{\dimexpr5cm+2\tabcolsep+\arrayrulewidth\relax}|}{\textcolor{ifscgreen}{\textbf{Semestre:}} \textbf{5º e 6º (História II)}} \\ \cline{2-3}
 & \textcolor{ifscgreen}{\textbf{CH EaD*:}} & \textcolor{ifscgreen}{\textbf{CH Total*:}} \\ \cline{2-3}
 & \textbf{00 h} & \textbf{80 h} \\ \hline
\endhead
\multicolumn{3}{|p{\dimexpr\linewidth-2\tabcolsep-2\arrayrulewidth\relax}|}{\textcolor{ifscgreen}{\textbf{Objetivos:}}} \\ \hline
\multicolumn{3}{|p{\dimexpr\linewidth-2\tabcolsep-2\arrayrulewidth\relax}|}{
\begin{itemize}[leftmargin=*,noitemsep,topsep=2pt]
 \item Analisar os sistemas econômicos ao longo da história do Brasil, destacando suas características, dinâmicas e impactos na organização social.
 \item Compreender o impacto das transformações econômicas na vida cotidiana brasileira, refletindo sobre as desigualdades socioeconômicas.
 \item Investigar as resistências e os projetos alternativos de organização socioeconômica (cooperativismo, economia solidária e outras formas de autogestão).
 \item Compreender o surgimento e a evolução dos direitos humanos, analisando os principais movimentos sociais que contribuíram para sua consolidação.
 \item Identificar os movimentos sociais brasileiros responsáveis pelo avanço na garantia de direitos para população negra, indígena, mulheres, LGBTQIA+, pessoas com deficiência e demais grupos vulneráveis.
 \item Analisar as relações entre sociedade e meio ambiente ao longo da história e os desafios para a sustentabilidade.
 \item Promover a conscientização sobre a importância da preservação ambiental e da adoção de práticas sustentáveis.
 \item Compreender os processos históricos de formação e transformação das identidades individuais e coletivas através da cultura.
 \item Analisar criticamente as expressões culturais e manifestações artísticas dos povos ao longo da história.
 \item Refletir sobre os impactos da globalização nas culturas locais e nos processos de hibridização e resistência cultural.
 \item Reconhecer e valorizar o papel das culturas africanas, afro-brasileiras e indígenas na formação da identidade social brasileira.
 \item Analisar criticamente as construções sociais das identidades de gênero e sexualidade ao longo da história.
\end{itemize}
} \\ \hline
\multicolumn{3}{|p{\dimexpr\linewidth-2\tabcolsep-2\arrayrulewidth\relax}|}{\textcolor{ifscgreen}{\textbf{Conteúdos:}}} \\ \hline
\multicolumn{3}{|p{\dimexpr\linewidth-2\tabcolsep-2\arrayrulewidth\relax}|}{
\begin{itemize}[leftmargin=*,noitemsep,topsep=2pt]
 \item Economia e Sociedade: sistemas econômicos ao longo do tempo e do espaço; impacto das transformações econômicas na vida das pessoas; mundos do trabalho na história; resistências e projetos alternativos de organização socioeconômica.
 \item Direitos Humanos e Justiça Social: história dos direitos contemporâneos (civis, políticos, socioeconômicos e humanos); movimentos sociais e lutas por justiça e igualdade no Brasil.
 \item Meio Ambiente e Sustentabilidade: relações entre sociedade e meio ambiente; impactos das atividades humanas no ecossistema global; capitalismo e a questão ambiental; movimentos ambientais.
 \item Ciência, Tecnologia e Inovação: avanços científicos e tecnológicos e seu impacto na sociedade; mudanças nas formas de produção e comunicação.
 \item Cultura e Identidade: patrimônio cultural; expressões culturais e manifestações artísticas; globalização e culturas locais; apropriação cultural; movimentos culturais marginalizados; culturas africanas, afro-brasileiras e indígenas; cultura popular e mídia.
 \item Identidades de Gênero e Sexualidade: elementos construídos sócio-historicamente.
\end{itemize}
} \\ \hline
\multicolumn{3}{|p{\dimexpr\linewidth-2\tabcolsep-2\arrayrulewidth\relax}|}{\textcolor{ifscgreen}{\textbf{Estratégias de Ensino e Aprendizagem:}}} \\ \hline
\multicolumn{3}{|p{\dimexpr\linewidth-2\tabcolsep-2\arrayrulewidth\relax}|}{- **Metodologia:** Abordagem participativa, crítica e contextualizada, amparada na pedagogia histórico-crítica, priorizando a relação horizontal entre professor e estudantes. Debates, análises de fontes históricas, atividades práticas e trabalhos em grupo, relacionando os conteúdos às experiências dos estudantes e à realidade social. Contextualização dos conteúdos históricos com acontecimentos e processos sociais atuais. Problematização das relações de poder e das desigualdades sociais na construção e interpretação da história.
- **Avaliação:** Não especificado no documento fonte.} \\ \hline
\multicolumn{3}{|p{\dimexpr\linewidth-2\tabcolsep-2\arrayrulewidth\relax}|}{\textcolor{ifscgreen}{\textbf{Bibliografia Básica:}}} \\ \hline
\multicolumn{3}{|p{\dimexpr\linewidth-2\tabcolsep-2\arrayrulewidth\relax}|}{1. BOULOS JÚNIOR, Alfredo. \textbf{360º: história sociedade \& cidadania}. São Paulo: FTD, 2015.
2. VICENTINO, Cláudio; DORIGO, Gianpaolo. \textbf{História geral e do Brasil}. São Paulo: Scipione, 2010.} \\ \hline
\multicolumn{3}{|p{\dimexpr\linewidth-2\tabcolsep-2\arrayrulewidth\relax}|}{\textcolor{ifscgreen}{\textbf{Bibliografia Complementar:}}} \\ \hline
\multicolumn{3}{|p{\dimexpr\linewidth-2\tabcolsep-2\arrayrulewidth\relax}|}{1. BHABHA, Homi. \textbf{O local da cultura}. Belo Horizonte: Editora UFMG, 1999.
2. CERTEAU, Michel de. \textbf{A invenção do cotidiano: artes de fazer}. Petrópolis, RJ: Vozes, 2009.
3. CAREGNATO, Célia Elizabete; BOMBASSARO, Luiz Carlos (org.). \textbf{Diversidade cultural: viver diferenças e enfrentar desigualdades na educação}. Erechim, RS: Novello \& Carbonelli, 2013.
4. GALEANO, Eduardo. \textbf{As caras e as máscaras}. Porto Alegre: L\&PM Editores, 2004.
5. HALL, Stuart. \textbf{A identidade cultural na pós-modernidade}. Rio de Janeiro: Lamparina, 2015.
6. HALL, Stuart; WOODWARD, Kathryn. \textbf{Identidade e diferença: a perspectiva dos estudos culturais}. Rio de Janeiro: Vozes, 2014.
7. HOBSBAWM, E. J. \textbf{Era dos extremos: o breve século XX: 1914-1991}. São Paulo: Companhia das Letras, 2005.
8. HOLLANDA, Heloísa Buarque de (org.). \textbf{Pensamento Feminista Hoje: perspectivas decoloniais}. Rio de Janeiro: Bazar do Tempo, 2020.
9. MIGNOLO, Walter. \textbf{Histórias Globais/projetos Locais}. Belo Horizonte: Editora UFMG, 2003.
10. SANTOS, Boaventura de Sousa; MARTINS, Bruno Sena (org.). \textbf{O pluriverso dos direitos humanos: a diversidade das lutas pela dignidade}. Belo Horizonte: Autêntica, 2019.
11. SOUZA, Laura Olivieri Carneiro de. \textbf{Quilombos: identidade e história}. Rio de Janeiro: Nova Fronteira, 2012.
12. RICKLEFS, Robert E. \textbf{A economia da natureza}. Rio de Janeiro: Guanabara Koogan, 2010.
13. WALSH, C. \textbf{Pedagogías decoloniales: prácticas insurgentes de resistir, (re)existir y (re)vivir. Tomo II}. Quito: Abya-Yala, 2017.} \\ \hline
\end{xltabular}

\vspace{0.4cm}


% ---------------------------------------------------------
% TABELA EMENTA 29: Sociologia — Ano 1
% ---------------------------------------------------------
\begin{xltabular}{\linewidth}{|X|p{2.5cm}|p{2.5cm}|}
\hline
\multirow{3}{=}{\textcolor{ifscgreen}{\textbf{Unidade Curricular:}}\\[3pt]\textbf{\large Sociologia — Ano 1}} & \multicolumn{2}{p{\dimexpr5cm+2\tabcolsep+\arrayrulewidth\relax}|}{\textcolor{ifscgreen}{\textbf{Semestre:}} \textbf{1º e 2º}} \\ \cline{2-3}
 & \textcolor{ifscgreen}{\textbf{CH EaD*:}} & \textcolor{ifscgreen}{\textbf{CH Total*:}} \\ \cline{2-3}
 & \textbf{00 h} & \textbf{80 h} \\ \hline
\endhead
\multicolumn{3}{|p{\dimexpr\linewidth-2\tabcolsep-2\arrayrulewidth\relax}|}{\textcolor{ifscgreen}{\textbf{Objetivos:}}} \\ \hline
\multicolumn{3}{|p{\dimexpr\linewidth-2\tabcolsep-2\arrayrulewidth\relax}|}{
\begin{itemize}[leftmargin=*,noitemsep,topsep=2pt]
 \item Analisar os diferentes discursos sobre a realidade: as explicações das Ciências Sociais amparadas nos vários paradigmas teóricos.
 \item Compreender as transformações que ocorrem nas sociedades humanas.
 \item Evidenciar a relação entre as questões individuais e as questões sociais.
 \item Formular questionamentos que permitam alcançar um conhecimento mais preciso da sociedade e uma postura crítica em relação às vivências que nos condicionam e limitam.
 \item Entender o processo de constituição, consolidação e desenvolvimento das sociedades modernas.
\end{itemize}
} \\ \hline
\multicolumn{3}{|p{\dimexpr\linewidth-2\tabcolsep-2\arrayrulewidth\relax}|}{\textcolor{ifscgreen}{\textbf{Conteúdos:}}} \\ \hline
\multicolumn{3}{|p{\dimexpr\linewidth-2\tabcolsep-2\arrayrulewidth\relax}|}{
\begin{itemize}[leftmargin=*,noitemsep,topsep=2pt]
 \item As Ciências Sociais e o objeto de estudo da Sociologia.
 \item Contexto de surgimento da Sociologia.
 \item Relação indivíduo/sociedade.
 \item Socialização.
 \item Instituições Sociais.
 \item Clássicos das Ciências Sociais.
 \item Cidadania e Direitos Humanos.
 \item Cultura e Ideologia.
 \item Cultura do ponto de vista antropológico.
 \item Multiculturalismo.
 \item Identidade e diferenças culturais.
 \item Cultura dominante e Indústria cultural.
 \item Subcultura/"Tribos urbanas".
\end{itemize}
} \\ \hline
\multicolumn{3}{|p{\dimexpr\linewidth-2\tabcolsep-2\arrayrulewidth\relax}|}{\textcolor{ifscgreen}{\textbf{Estratégias de Ensino e Aprendizagem:}}} \\ \hline
\multicolumn{3}{|p{\dimexpr\linewidth-2\tabcolsep-2\arrayrulewidth\relax}|}{- **Metodologia:** Ensino desenvolvido através da leitura, análise, discussão e exposição de textos e imagens em sala de aula, visando o exercício do debate e da reflexão crítica; aulas expositivas e dialogadas; recursos didático-pedagógicos como filmes, seminários, documentários e entrevistas; estímulo à autonomia investigativa. Recursos: textos (livros, apostilas, artigos, estudos de caso), quadro e pincel, recursos audiovisuais, equipamentos de informática.
- **Avaliação:** Não especificado no documento fonte.} \\ \hline
\multicolumn{3}{|p{\dimexpr\linewidth-2\tabcolsep-2\arrayrulewidth\relax}|}{\textcolor{ifscgreen}{\textbf{Bibliografia Básica:}}} \\ \hline
\multicolumn{3}{|p{\dimexpr\linewidth-2\tabcolsep-2\arrayrulewidth\relax}|}{1. Livro didático fornecido pelo Fundo Nacional de Desenvolvimento da Educação (FNDE).
2. OLIVEIRA, Pérsio Santos de. \textbf{Introdução à Sociologia: Série Brasil}. São Paulo: Editora Ática, 2011.} \\ \hline
\multicolumn{3}{|p{\dimexpr\linewidth-2\tabcolsep-2\arrayrulewidth\relax}|}{\textcolor{ifscgreen}{\textbf{Bibliografia Complementar:}}} \\ \hline
\multicolumn{3}{|p{\dimexpr\linewidth-2\tabcolsep-2\arrayrulewidth\relax}|}{1. COSTA, Maria Cristina. \textbf{Sociologia: introdução à ciência da sociedade}. 5. ed. São Paulo: Editora Moderna, 2016.
2. OLIVEIRA, Luiz Fernandes de; COSTA, Ricardo Cesar Rocha da. \textbf{Sociologia para jovens do século XXI}. Rio de Janeiro: Imperial Novo Milênio, 2013.
3. HELENA, Bomeny et al. \textbf{Tempos Modernos, tempos de sociologia}. 4. ed. São Paulo: Ed. do Brasil, 2016.} \\ \hline
\end{xltabular}

\vspace{0.4cm}


% ---------------------------------------------------------
% TABELA EMENTA 30: Sociologia — Ano 2
% ---------------------------------------------------------
\begin{xltabular}{\linewidth}{|X|p{2.5cm}|p{2.5cm}|}
\hline
\multirow{3}{=}{\textcolor{ifscgreen}{\textbf{Unidade Curricular:}}\\[3pt]\textbf{\large Sociologia — Ano 2}} & \multicolumn{2}{p{\dimexpr5cm+2\tabcolsep+\arrayrulewidth\relax}|}{\textcolor{ifscgreen}{\textbf{Semestre:}} \textbf{3º e 4º}} \\ \cline{2-3}
 & \textcolor{ifscgreen}{\textbf{CH EaD*:}} & \textcolor{ifscgreen}{\textbf{CH Total*:}} \\ \cline{2-3}
 & \textbf{00 h} & \textbf{80 h} \\ \hline
\endhead
\multicolumn{3}{|p{\dimexpr\linewidth-2\tabcolsep-2\arrayrulewidth\relax}|}{\textcolor{ifscgreen}{\textbf{Objetivos:}}} \\ \hline
\multicolumn{3}{|p{\dimexpr\linewidth-2\tabcolsep-2\arrayrulewidth\relax}|}{
\begin{itemize}[leftmargin=*,noitemsep,topsep=2pt]
 \item Analisar os diferentes discursos sobre a realidade: as explicações das Ciências Sociais amparadas nos vários paradigmas teóricos.
 \item Compreender as transformações que ocorrem nas sociedades humanas.
 \item Evidenciar a relação entre as questões individuais e as questões sociais.
 \item Formular questionamentos que permitam alcançar um conhecimento mais preciso da sociedade e uma postura crítica em relação às vivências que nos condicionam e limitam.
 \item Entender o processo de constituição, consolidação e desenvolvimento das sociedades modernas.
\end{itemize}
} \\ \hline
\multicolumn{3}{|p{\dimexpr\linewidth-2\tabcolsep-2\arrayrulewidth\relax}|}{\textcolor{ifscgreen}{\textbf{Conteúdos:}}} \\ \hline
\multicolumn{3}{|p{\dimexpr\linewidth-2\tabcolsep-2\arrayrulewidth\relax}|}{
\begin{itemize}[leftmargin=*,noitemsep,topsep=2pt]
 \item Política e Cidadania — Etimologia, História e conceituação das palavras.
 \item Clássicos da Política.
 \item Estado, Governo e Nação.
 \item Formas de organização do Estado.
 \item Regimes de Governo.
 \item Estratificação Social.
 \item Classes sociais e desigualdades; Castas e Estamentos.
 \item Movimentos sociais, sindicalismo e democracia.
 \item Novos Movimentos Sociais.
 \item Cidadania, Direitos Humanos e Meio Ambiente.
\end{itemize}
} \\ \hline
\multicolumn{3}{|p{\dimexpr\linewidth-2\tabcolsep-2\arrayrulewidth\relax}|}{\textcolor{ifscgreen}{\textbf{Estratégias de Ensino e Aprendizagem:}}} \\ \hline
\multicolumn{3}{|p{\dimexpr\linewidth-2\tabcolsep-2\arrayrulewidth\relax}|}{- **Metodologia:** Ensino desenvolvido através da leitura, análise, discussão e exposição de textos e imagens em sala de aula, visando o exercício do debate e da reflexão crítica; aulas expositivas e dialogadas; recursos didático-pedagógicos como filmes, seminários, documentários e entrevistas; estímulo à autonomia investigativa. Recursos: textos (livros, apostilas, artigos, estudos de caso), quadro e pincel, recursos audiovisuais, equipamentos de informática.
- **Avaliação:** Não especificado no documento fonte.} \\ \hline
\multicolumn{3}{|p{\dimexpr\linewidth-2\tabcolsep-2\arrayrulewidth\relax}|}{\textcolor{ifscgreen}{\textbf{Bibliografia Básica:}}} \\ \hline
\multicolumn{3}{|p{\dimexpr\linewidth-2\tabcolsep-2\arrayrulewidth\relax}|}{1. Livro didático fornecido pelo Fundo Nacional de Desenvolvimento da Educação (FNDE).
2. OLIVEIRA, Pérsio Santos de. \textbf{Introdução à Sociologia: Série Brasil}. São Paulo: Editora Ática, 2011.} \\ \hline
\multicolumn{3}{|p{\dimexpr\linewidth-2\tabcolsep-2\arrayrulewidth\relax}|}{\textcolor{ifscgreen}{\textbf{Bibliografia Complementar:}}} \\ \hline
\multicolumn{3}{|p{\dimexpr\linewidth-2\tabcolsep-2\arrayrulewidth\relax}|}{1. COSTA, Maria Cristina. \textbf{Sociologia: introdução à ciência da sociedade}. 5. ed. São Paulo: Editora Moderna, 2016.
2. OLIVEIRA, Luiz Fernandes de; COSTA, Ricardo Cesar Rocha da. \textbf{Sociologia para jovens do século XXI}. Rio de Janeiro: Imperial Novo Milênio, 2013.
3. HELENA, Bomeny et al. \textbf{Tempos Modernos, tempos de sociologia}. 4. ed. São Paulo: Ed. do Brasil, 2016.} \\ \hline
\end{xltabular}

\vspace{0.4cm}

